\begin{abstract}
Despite the widespread adoption of Transformer models for NLP tasks, the expressive power of these models is not well-understood. In this paper, we establish that Transformer models are universal approximators of continuous {\em permutation equivariant} sequence-to-sequence functions with compact support, which is quite surprising given the amount of shared parameters in these models. Furthermore, using positional encodings, we circumvent the restriction of permutation equivariance, and show that Transformer models can universally approximate {\em arbitrary} continuous sequence-to-sequence functions on a compact domain. Interestingly, our proof techniques clearly highlight the different roles of the self-attention and the feed-forward layers in Transformers. In particular, we prove that fixed width self-attention layers can compute \emph{contextual mappings} of the input sequences, playing a key role in the universal approximation property of Transformers. Based on this insight from our analysis, we consider other simpler alternatives to self-attention layers and empirically evaluate them.
\end{abstract}


%%%%%%%%%%%%%%%%%%%%%%%%%%%%%%%%%%%%%%%%%%%%%%%%%%%%%%%%%%%%%%%%%%%%%%%%%%%%%%%%%%%%%%%%%%%%%%%%%%%

%Despite the widespread adoption of Transformer models for NLP tasks, the expressive power of these models is not well understood. In this paper, we establish that Transformer models are universal approximators of {\em permutation equivariant} sequence-to-sequence functions with compact support, which is quite surprising given the amount of shared parameters in these models. Interestingly, our proof techniques clearly highlight the different roles of the  self-attention and feed-forward layers in Transformers. We prove that fixed width self-attention layers can compute contextual mappings for input of any length, playing a key role in the universal approximation property of Transformers. Based on this insight from our analysis, we empirically evaluate other architectures that can compute contextual mappings and compare them with Transformers. \todo{Should we say something stronger here?}

%%%%%%%%%%%%%%%%%%%%%%%%%%%%%%%%%%%%%%%%%%%%%%%%%%%%%%%%%%%%%%%%%%%%%%%%%%%%%%%%%%%%%%%%%%%%%%%%%%

% In their seminal work, \citet{vaswani2017attention} demonstrated the power of attention mechanisms by proposing a novel architecture with state-of-the-art performance, namely {\em Transformer} model, that completely does away with recurrence and convolution blocks which are integral components of conventional sequence-to-sequence models. Since then, Transformer models have become increasingly prevalent by leading a series of breakthroughs in the natural language understanding space~\citep{devlin2018bert, radford2018gpt, radford2019gpt2, xlnet2019, roberta2019}.

% Despite the widespread adoption of Transformer models, many basic properties of these models are far from well-understood. This paper explores the expressive power of Transformer models as sequence to sequence functions. We establish that Transformer models are universal approximators of {\em permutation equivariant} sequence-to-sequence functions with compact support, which is quite surprising given the amount of shared parameters in these models. Furthermore, we rigorously show that employing positional embeddings significantly enhances the expressive power of Transformer models beyond permutation equivariant functions as they become universal approximators of all sequence-to-sequence functions on a compact domain. Interestingly, our proof techniques clearly highlight the different roles of two key components of Transformer models, namely self-attention and feed-forward layers. In particular, the self-attention layers and the feed-forward layers help the model generate contextual embeddings and {\color{red}create functional mapping/memorization, respectively. Based on this insight from our analysis, we propose novel architectural modifications to Transformer models. We further demonstrate the utility of theses architectural changes through our empirical evaluation.}
